\documentclass[12pt, a4paper]{article}

% 中文支持
\usepackage{ctex}
\usepackage{geometry}
\usepackage{hyperref}
\usepackage{booktabs}
\usepackage{graphicx}
\usepackage{enumitem}
\usepackage{titlesec}
\usepackage{parskip}

% 页面设置
\geometry{left=2.5cm, right=2.5cm, top=2.5cm, bottom=2.5cm}

% 标题格式
\titleformat{\section}{\bfseries\large}{}{0em}{}
\titleformat{\subsection}{\bfseries\normalsize}{}{0em}{}

% 超链接设置
\hypersetup{
    colorlinks=true,
    linkcolor=blue,
    urlcolor=blue
}

\begin{document}

% ========== 封面信息 ==========
\begin{center}
    {\Large \textbf{2026年第十三届中国可视化与可视分析大会}}\\[0.5em]
    {\large \textbf{数据可视化竞赛 赛道1-II}}\\[0.3em]
    {\large \textbf{（ChinaVis Data Challenge 2026 - mini challenge 1-II）}}\\[1em]
    {\LARGE \textbf{答\quad 卷}}
\end{center}

\vspace{1em}
\noindent\textbf{参赛队名称：}天津大学-张三-赛道1-I

\vspace{0.5em}
\noindent\textbf{团队成员：}

\begin{tabbing}
    \hspace{4em} \= 张三，天津大学，zs@tju.edu.cn，队长 \\
    \> 李四，天津大学，\href{mailto:ls@tju.edu.cn}{ls@tju.edu.cn} \\
    \> 黄五，天津大学，hw@tju.edu.cn，指导老师
\end{tabbing}

\noindent\textbf{团队成员是否与报名表一致（是或否）：}是

\noindent\textbf{是否学生队（是或否）：}是

\noindent\textbf{使用的分析工具或开发工具}（如果使用了自己研发的软件或工具请具体说明）：\\
D3，Excel，MySQL，Qt，Deepseek，CVASTer（天津大学xxx中心开发的数据可视分析工具）

\noindent\textbf{共计耗费时间（人天）：}60人天

\noindent\textbf{本次比赛结束后，我们是否可以在网络上公布该答卷与视频（是或否）：}是

\vspace{1em}
\hrule
\vspace{1.5em}

% ========== 题目1 ==========
\section*{1、采用体数据渲染技术，并设计合适的传递函数和光照效果，展示宇宙演化中密度信息的变化。}

{\small （建议参赛者回答此题文字不多于800字，图片不多于5张）}

\vspace{0.5em}
\noindent（下面是答题区域）

\noindent 本题采用对数归一化和发光点云近似体渲染方式展示 Nyx 密度场的时序演化。原始数据以 little-endian float32 读取，并按照 $z,y,x$ 顺序恢复为 $128^3$ 三维体数据。由于宇宙密度值动态范围较大，直接线性映射会导致低密度背景和高密度节点难以同时观察，因此先对密度取 $\log_{10}$，再使用全局分位数范围进行归一化。

\noindent 传递函数设计为分层颜色和透明度映射：低到中等密度区域使用蓝紫色半透明发光点表示宇宙网的连续云雾和丝状桥接结构；高密度区域使用玫紫色增强亮度；极高密度核心使用金白色强发光材质突出节点和团块。渲染中加入透明参考框、深空背景、面积光、轮廓光和泛光效果，使三维结构的空间层次更清晰。

\noindent 从视频结果看，早期密度分布相对分散，发光结构弱且连接关系不明显。随着时间推进，高密度区域逐渐沿纤维状通道汇聚，节点亮度和尺度增强，宇宙网骨架变得连续清晰。该可视化方式能够同时表达密度强弱、空间位置和时间变化，适合观察宇宙密度场由近均匀状态向“纤维--节点”结构演化的过程。

% ========== 题目2 ==========
\section*{2、基于可视化图像结果，归纳宇宙密度的演化特征，总结宇宙结构形成、星系际介质演化的核心规律，阐释可视化技术在宇宙学数据分析中的应用价值。}

{\small （建议参赛者回答此题文字不多于800字，图片不多于5张）}

\vspace{0.5em}
\noindent（下面是答题区域）

\noindent 本题采用统一对数密度标尺对多个关键时间步进行可视化比较。首先读取 little-endian float32 体数据并按 $z,y,x$ 顺序恢复为三维密度场，再通过对数归一化压缩密度动态范围，突出高密度骨架与低密度空洞之间的差异。视频中使用固定候选点云和 Shape Key 方式保持空间拓扑一致，使不同时刻的亮度、半径和连接关系可以连续比较。

\noindent 从可视化结果看，早期密度分布相对平滑，高密度区域主要表现为离散斑块；随着演化推进，物质逐渐沿引力势阱汇聚，离散斑块被拉伸、连接为丝状纤维，并在纤维交汇处形成更亮、更稳定的节点。低密度区域则被进一步排空，形成更清晰的宇宙空洞。整体演化呈现“空洞--纤维--节点”的层级宇宙网结构。

\noindent 这说明星系际介质并非均匀扩散，而是在暗物质引力作用下沿宇宙网骨架输运和压缩，高密度节点可视为星系团或强吸收结构的候选区域。可视化技术的价值在于把三维空间结构和时间演化过程直观结合，帮助快速识别结构形成阶段、定位致密节点和纤维通道，并为后续定量统计和物理解释提供依据。

% ========== 题目3 ==========
\section*{3、宇宙密度时序统计特征分析：早期宇宙物质分布相对均匀，密度数值高度集中于均值区间。在100个演化时间步长的暗物质引力持续作用下，宇宙物质不断聚集坍缩，整体呈现显著的"团块化"特征，密度数值波动幅度持续增大，逐步形成极端低密度宇宙空洞与极端高密度结构峰值的两极分化分布。自主统计每个时间步长的气体密度数据，构建密度对数直方图，量化追踪全域密度分布随时间的演化偏移规律。}

{\small （建议参赛者回答此题文字不多于800字，图片不多于5张）}

\vspace{0.5em}
\noindent（下面是答题区域）

\noindent 本题对 100 个时间步的气体密度数据进行统一统计分析。每个时间步先按 little-endian float32 读取并恢复为 $128^3$ 体数据，再进行采样和 $\log_{10}$ 变换，构建密度对数直方图。同时计算均值、标准差、p99、p99.9 和全局 Top 1\% 阈值占比，用于追踪分布宽度、高密度尾部和极端峰值随时间的变化。

\noindent 统计动画显示，早期密度直方图集中在较窄区间内，说明物质分布相对均匀。随着演化推进，直方图逐渐展宽，低密度端和高密度端同时增强：一方面大量区域密度降低，对应宇宙空洞的形成；另一方面少量体素进入高密度尾部，对应宇宙网节点和致密纤维的增强。标准差持续增大反映密度波动幅度扩大，p99 与 p99.9 上升则说明高密度尾部不断增长。

\noindent 因此，Nyx 密度场在时间上呈现由均匀到团块化、由单峰集中到两极分化的演化规律。密度对数直方图提供了全域分布视角，高分位曲线则把极端高密度结构的增长量化出来，两者共同证明暗物质引力持续驱动物质聚集坍缩，并逐步形成极端低密度空洞与极端高密度峰值并存的宇宙大尺度结构。

% ========== 题目4 ==========
\section*{4、相空间交互式刷选可视分析：搭建联动式交互可视分析仪表盘，实现基于相空间的筛选分析功能。支持用户在统计直方图中自定义框选特定密度区间，例如筛选占比1\%的极高密度区间。系统可实时联动三维空间视图，快速渲染匹配该密度阈值的物理单元格数据，直观验证直方图高密度尾部数据与宇宙网节点致密结构的对应关系，实现统计特征与空间物理结构的双向关联分析。}

{\small （建议参赛者回答此题文字不多于800字，图片不多于5张）}

\vspace{0.5em}
\noindent（下面是答题区域）

\noindent 本题构建相空间联动筛选仪表盘，将统计直方图中的密度区间选择与三维空间视图中的物理单元格高亮绑定。系统先对关键时间步的密度数据进行 $\log_{10}$ 变换并构建全局密度范围，再在左侧直方图中显示当前时间步的密度分布。黄色阈值窗口表示被筛选的高密度区间，示例中从高密度尾部逐步移动到约 Top 1\% 的极端密度范围。

\noindent 右侧三维空间视图使用与阈值区间匹配的候选体素点云进行高亮显示。为了保证 Blender 中动画播放稳定，脚本将多个关键时间步的筛选结果预先烘焙为 Shape Key：当直方图阈值窗口变化时，空间视图中对应的高密度体素同步增大、发光或收缩，从而模拟交互刷选时统计视图和空间视图之间的联动关系。

\noindent 可视分析结果表明，直方图右侧高密度尾部对应的体素并非随机散布，而是集中出现在宇宙网节点、纤维交汇处和局部致密团块区域。当筛选区间收紧到极高密度范围时，高亮区域明显收缩到少数核心节点。这验证了统计特征与空间物理结构之间的对应关系，也说明联动刷选能够帮助用户从全局密度分布快速定位物理上重要的宇宙网结构。

\end{document}
% ========== 附录 ==========
\section*{附录：数据与脚本说明}
本作品所有原始数据存放于 \texttt{Nyx} 文件夹，所有可视化脚本与网页展示文件均在主目录。请确保解压后文件结构如下：

\begin{itemize}
    \item Nyx/（原始数据，100个dat文件）
    \item NYX01\_FIXED\_RUN\_THIS.py
    \item NYX02\_FIXED\_RUN\_THIS.py
    \item NYX03\_FIXED\_RUN\_THIS.py
    \item nyx\_04\_interactive\_linked\_dashboard.py
    \item index.html
    \item README.md
\end{itemize}

如脚本因路径问题无法运行，请参考 README.md 修改脚本中的数据路径为本地实际路径。
